\begin{lemma}
\label{editorial-source-lemma-open-going-down}
Let $R \to S$ be a ring map. If the induced map
$\varphi : \Spec(S) \to \Spec(R)$ is open, then
$R \to S$ satisfies going down.
\end{lemma}

\begin{proof}
Retain $\mathfrak p\subseteq\mathfrak p'$ and the original
$\mathfrak q'$ lying over $\mathfrak p'$.
For $g\notin\mathfrak q'$, the open set $D(g)$ contains
$\mathfrak q'$, so its image is an open neighbourhood of
$\mathfrak p'$. This image contains a standard open
$D(f)$ with $f\notin\mathfrak p'$. Then
$f\notin\mathfrak p$, so it also contains $\mathfrak p$.
The fibre criterion in Lemma \ref{editorial-source-lemma-in-image} gives
$S_g\otimes_R\kappa(\mathfrak p)\ne0$ for every such $g$.

Here is the precise directed localization argument. Index
finite subsets $E\subseteq S\setminus\mathfrak q'$ by
inclusion and set $g_E=\prod_{g\in E}g$, with $g_\varnothing=1$.
For $E\subseteq F$ the map is
$$
S_{g_E}\longrightarrow S_{g_F},\qquad
\frac{s}{g_E^n}\longmapsto
\frac{s\,g_{F\setminus E}^{\,n}}{g_F^n}.
$$
It is the original localization map, since all factors
in the product $g_F$ become units. The map from their
colimit to $S_{\mathfrak q'}$ sends each fraction to the
same original fraction. It is surjective: $s/t$ comes
from $E=\{t\}$. If $s/g_E^n$ maps to zero, there is
$u\notin\mathfrak q'$ with $us=0$. At the stage
$F=E\cup\{u\}$, $u$ is a unit, so the image of $s$
and then that fraction is zero. This proves injectivity
and the representative formula for the inverse. It is
the comparison of Lemma \ref{algebra-lemma-localization-colimit}.

Tensor-colimit compatibility gives the unital ring
comparison
$$
S_{\mathfrak q'}\otimes_R\kappa(\mathfrak p)
\cong\colim_E(S_{g_E}\otimes_R\kappa(\mathfrak p)),
$$
sending a stage tensor $s/g_E^n\otimes c$ to the tensor
of the same fraction and $c$; the inverse follows by
representing the finitely many fractions at a common
stage, as in Lemma \ref{editorial-source-lemma-tensor-products-commute-with-limits}.
Every stage ring is nonzero by the previous paragraph,
because $g_E\notin\mathfrak q'$. Their colimit is nonzero:
if its unit were zero, it would become zero at a later
stage, but every transition is unital and that later
ring is nonzero.

Choose a prime of this nonzero fibre ring and pull it
back along $S\to S_{\mathfrak q'}\to
S_{\mathfrak q'}\otimes_R\kappa(\mathfrak p)$ to a prime
$\mathfrak q$ of $S$. Every element outside
$\mathfrak q'$ maps to a unit, so
$\mathfrak q\subseteq\mathfrak q'$.
The contraction to the field $\kappa(\mathfrak p)$ is
its zero prime, and the original map $R\to\kappa(\mathfrak p)$
has kernel $\mathfrak p$. Thus the contraction of
$\mathfrak q$ to $R$ is exactly $\mathfrak p$.

In fact this argument proves the exact weaker criterion:
$R\to S$ satisfies going down if and only if the image
of every principal open $D(g)\subseteq\Spec(S)$ is stable
under generalization. For sufficiency that stability
supplies $\mathfrak p$ in each image containing
$\mathfrak p'$, and the remaining argument just given
applies without openness. Conversely, if going down
holds and $\mathfrak q'\in D(g)$ lies over
$\mathfrak p'$, every $\mathfrak p\subseteq\mathfrak p'$
has a lift $\mathfrak q\subseteq\mathfrak q'$.
Since $g\notin\mathfrak q'$, also $g\notin\mathfrak q$,
so the lift still belongs to the original $D(g)$.
This proves necessity and both directions of the criterion.

Openness is stronger than this criterion. The flat map
$\mathbf Z\to\mathbf Q$ satisfies going down by
Lemma \ref{editorial-source-lemma-flat-going-down}, but the image of its
whole spectrum is $\{(0)\}$, which is not open in
$\Spec(\mathbf Z)$. Indeed any open neighbourhood of
$(0)$ contains $D(n)$ for some nonzero integer $n$.
A prime divisor $p$ of $|n|+1>1$ does not divide $n$,
so $(p)\in D(n)$ and $(p)\ne(0)$. This proves the
failure of openness. Flatness of $\mathbf Q$ was proved
in Example \ref{editorial-source-example-flat-infinite-intersections}.
\end{proof}